\subsection{Maximal test rank in a nonabelian variety: Problem 16.87}
\label{cand:16.87}

A test set in a group is a set whose pointwise stabilizer among
endomorphisms consists of automorphisms. The test rank is the least
size of such a set. Both parts of the question have a negative answer
for the variety
\[
 \mathcal M_p=\{G:[[x,y],z]=1,\ [x,y]^p=1\},
\]
where \(p\) is any prime. It is nonabelian because it contains
\(\operatorname{UT}_3(\F_p)\), and nonperiodic because it contains
\(\Z\).

\begin{proposition}
The relatively free group \(G_r\) of rank \(r\geq1\) in
\(\mathcal M_p\) has test rank \(r\).
\end{proposition}
\begin{proof}
Its abelianization is \(\Z^r\): the generator maps to and from
\(\Z^r\) give inverse maps on abelianization, since \(\Z^r\)
belongs to the variety. Denote the resulting map by
\(v:G_r\to\Z^r\). If \(S\) has fewer than \(r\) elements,
choose a primitive nonzero integral row vector \(w\) annihilating
all \(v(s)\), \(s\in S\), by solving the corresponding homogeneous
linear equations over \(\Q\). Then \(f=wv:G_r\to\Z\) is onto
and vanishes on \(S\).

Each free generator \(x_i\) has central \(p\)-th power, because
\([x_i^p,g]=[x_i,g]^p=1\). Choose \(i\) with \(w_i\ne0\)
and put \(c=x_i^{p\,\operatorname{sign}(w_i)}\). The map
\[
 \phi(g)=g c^{f(g)}
\]
is an endomorphism by centrality of \(c\), and fixes \(S\)
pointwise. But
\(f(\phi(G_r))=(1+p|w_i|)\Z\) is proper in \(\Z\), so
\(\phi\) is not surjective. Thus no smaller set is a test set.
The \(r\) free generators themselves form one. For \(r=1\)
the same argument excludes the empty set.
\end{proof}

The free groups here have torsion in their commutator subgroups;
nonperiodicity of the variety does not mean that its free groups
are torsion-free. No such torsion-free hypothesis is printed in
the problem. The proof uses only central powers and integer
linear algebra. Source: \artifact{research/16.87-proof.md}.

